% !PRINTED 30
\weberpage{30}

\begingroup
\setlength{\abovedisplayskip}{5pt plus 2pt minus 1pt}
\setlength{\belowdisplayskip}{5pt plus 2pt minus 1pt}
\setlength{\abovedisplayshortskip}{3pt plus 1pt}
\setlength{\belowdisplayshortskip}{4pt plus 1pt}
\linespread{0.94}\selectfont

\begin{center}
{\large \S.\ 4.}

\vspace{0.5em}
{\large\gesperrt{Theilung durch eine lineare Function.}}
\end{center}

Wir wollen die allgemeine Vorschrift für die Division noch auf einen
besonderen Fall anwenden. Es sei der Dividend
\begin{equation}
 f(x)=a_0x^n+a_1x^{n-1}+a_2x^{n-2}+\cdots+a_n
\tag{1}
\end{equation}
beliebig, dagegen der Divisor vom ersten Grade oder, wie man auch sagt,
linear. Wir wollen auch den Coëfficienten der ersten Potenz von $x$ im
Divisor $=1$ voraussetzen und also den Divisor in der einfachen Form
$(x-\alpha)$ annehmen, worin $\alpha$ beliebig bleibt. Der Quotient
$Q$ ist in diesem Falle vom $(n-1)$\textsuperscript{ten} Grade und der
Rest $C$ vom $0$\textsuperscript{ten} Grade, d. h. von $x$ unabhängig.
Setzen wir also
\begin{equation}
 f(x)=(x-\alpha)Q+C,
\tag{2}
\end{equation}
so enthält $C$ das $x$ nicht mehr, und wenn wir
\begin{equation}
 Q=q_0x^{n-1}+q_1x^{n-2}+\cdots+q_{n-2}x+q_{n-1}
\tag{3}
\end{equation}
setzen, so folgt aus (2):
\begin{equation}
\begin{aligned}
 f(x)={}&q_0x^n+q_1x^{n-1}+\cdots+q_{n-2}x^2+q_{n-1}x+C\\
 &-\alpha q_0x^{n-1}-\cdots-\alpha q_{n-3}x^2
   -\alpha q_{n-2}x-\alpha q_{n-1},
\end{aligned}
\tag{4}
\end{equation}
und aus der Vergleichung mit (1):
\begin{equation}
\begin{aligned}
 q_0&=a_0,\\
 q_1-\alpha q_0&=a_1,\\
 q_2-\alpha q_1&=a_2,\\
 &\cdots,\\
 q_{n-1}-\alpha q_{n-2}&=a_{n-1},\\
 C-\alpha q_{n-1}&=a_n.
\end{aligned}
\tag{5}
\end{equation}
Daraus erhält man
\begin{equation}
\begin{aligned}
 q_0&=a_0,\\
 q_1&=a_0\alpha+a_1,\\
 q_2&=a_0\alpha^2+a_1\alpha+a_2,\\
 &\cdots,\\
 q_{n-1}&=a_0\alpha^{n-1}+a_1\alpha^{n-2}+\cdots+a_{n-1},\\
 C&=a_0\alpha^n+a_1\alpha^{n-1}+\cdots+a_{n-1}\alpha+a_n=f(\alpha).
\end{aligned}
\tag{6}
\end{equation}
$C$ entsteht aus $f(x)$, wenn man $x=\alpha$ setzt, und kann also auch
mit $f(\alpha)$ bezeichnet werden.
\endgroup
